% ============================================================
\section{Análise Estatística dos Resultados}
\label{sec:analise_estatistica}
% ============================================================

Dado o reduzido tamanho amostral, a diferença de desempenho entre o
\textit{XGBoost} e o SVM foi submetida a um teste de significância antes de se
afirmar a superioridade de um sobre o outro (BROWNLEE, 2020b). A comparação
usou os valores de AUC pareados nas dez partições da validação cruzada, de
modo que ambos os modelos foram avaliados sobre exatamente os mesmos
\textit{folds}.

O \textit{XGBoost} + SMOTE (AUC $= 0,722$) superou o SVM + SMOTE
(AUC $= 0,678$) por uma diferença média de $0,044$, vencendo em 6 dos 10
\textit{folds}. Ainda assim, o teste \textit{t} de Student pareado não rejeitou
a hipótese de igualdade ($t = 1,19$; $\text{gl} = 9$; abaixo do crítico
$2,262$), e o intervalo de confiança de 95\,\% da diferença ($-0,039$ a
$+0,127$) contém o zero. O teste não-paramétrico de Wilcoxon confirmou o
resultado ($W = 18,5$; $n = 10$). Ou seja, a AUC isolada não sustenta a
superioridade do \textit{XGBoost} de forma estatisticamente significativa,
achado atribuído à baixa potência imposta pelos poucos casos positivos.

A separação entre os modelos torna-se clara nas métricas de concordância. O
Coeficiente Kappa de Cohen do \textit{XGBoost} (0,372 sob \textit{BorderlineSMOTE})
supera o melhor resultado do SVM (0,210), padrão que se repete no F1-Score.
Ainda assim, na escala de Landis e Koch (LANDIS; KOCH, 1977), nenhuma
configuração atinge concordância moderada ($\kappa > 0,40$), permanecendo na
faixa razoável (\textit{fair}). Isso reforça a recomendação dos modelos como
ferramentas auxiliares de triagem, e não como substitutos da avaliação clínica.
